\section{Neutrino Energy Calibration\label{sec:CALIBRATION}}
Once the reconstruction algorithms are defined and we have obtained the raw variables from the TPC and PDS subsystems, we can proceed to calibrate the neutrino energy. The calibration process is based on the true information of the main electron of the interaction, which allows us to compute a hit-dependent calibration workflow that recovers the particle's original kinetic energy. The average calibration factor is then applied to all adjacent clusters of the event candidate and the total neutrino energy is computed from four different algorithms (see \cref{sec:NEUTRINO_CALIBRATION}). During the \gls{tpc} calibration process, the drift coordinate is assumed from truth. A discussion on how the intrinsic background contamination of the detector affects the signal reconstruction at these low energies and the effects on drift coordinate resolution (as a result of \gls{tpc}--\gls{pds} match) are presented in~\cref{sec:RESULTS}.

\subsection{Electron-Lifetime Correction}\label{sec:ELECTRON_LIFETIME}
\gls{tpc} hits are first grouped into spatial clusters using a dedicated clustering algorithm. Each event contains a \textbf{main cluster}, corresponding to the primary outgoing electron, and zero or more \textbf{adjacent clusters}, interpreted as gamma-induced ``blips''. 

In order to calibrate the charge attributed to each of these clusters, the first conceptual step is to understand that there is a direct correlation between the particle's energy and the collected charge. Keeping aside the differences in electron yield that each kind of particle produces and focusing on the main electron of the \gls{cc} $\nu-\text{Ar}$ interaction, we identify the main cluster as the one with the highest charge and assume that it corresponds to a fully-contained electron. This assumption allows us to use the known true energy of the main electron (from simulation) as a reference for calibrating the charge-to-energy conversion factor:

\begin{align}
    Q^{\text{coll}}(n_\text{hit}) = f\left(E^{\text{true}}_e,\,\overrightarrow{p}\,,n_\text{hit}\right) \cdot E^{\text{true}}_e\,.
\end{align}

Here, \( Q^{\text{coll}} \) is the charge collected in the main cluster, \( E^{\text{true}}_e \) is the true energy of the main electron, and \( f \) is the theoretical charge-to-energy conversion factor that depends on the true energy, the vertex of the electron's interaction (represented by \( \overrightarrow{p} \)), and the number of hits in the resulting cluster (\( n_\text{hit} \)). 

From an experimental perspective, the charge collected in the main cluster is affected by various factors, that are a priory unknown, including the particles energy and position. We define an effective charge-to-energy conversion factor, \( F_0(n_\text{hit})\), that encapsulates all these effects as a function of the number of hits in the cluster:

\begin{align}
    F_0(n_\text{hit}) = \dfrac{1}{\Delta t}\int^{\Delta t}_0\dfrac{Q^\text{corr}(t\,,n_\text{hit})}{E^\text{true}_e}dt \quad &\text{with} \quad Q^\text{corr}(n_\text{hit}) = \dfrac{Q^\text{coll}(n_\text{hit})}{e^{-\frac{t}{\tau_e}}}\,. \\
    E^\text{cluster}_e &= F_0(n_\text{hit}) \cdot Q^\text{corr}(n_\text{hit})\,.
\end{align}\label{eq:ELECTRON_LIFETIME}

Here \( Q^\text{corr} \) is the lifetime-corrected cluster charge, \( E^\text{cluster}_e \) is the computed electron energy for the cluster, \( \Delta t \) is the maximum drift time of the electron, and \( \tau_e \) is the electron lifetime measured here independently and cross-checked against the simulation input.

The first step in calibration involves converting the total collected charge of each cluster to a computed electron energy by correcting for the electron-lifetime attenuation in the \gls{lar} (see~\cref{eq:ELECTRON_LIFETIME}). \Cref{fig:DRIFT_ATTENUATION_CORRECTION} shows the result of this procedure.

\begin{figure}[ht!]
    \centering
    \includegraphics[width=.49\linewidth]{Images/Reco/charge_lifetime_correction_hd_1x2x6_marley_official_corrected_nhits_nan_hist1d.png}
    \includegraphics[width=.49\linewidth]{Images/Reco/charge_lifetime_correction_vd_1x8x14_3view_30deg_marley_official_corrected_nhits_nan_hist1d.png}
    \caption[width=1\linewidth]{Drift attenuation correction for \gls{hd} (left) and \gls{vd} (right).\label{fig:DRIFT_ATTENUATION_CORRECTION}}
\end{figure}

\Cref{fig:CHARGE_CORRECTION_FACTOR} shows the charge convertion factor as a function of the clusters $n_\text{hit}$. Due to the nature of clustering, the first points of the distribution are affected by threshold effects, which can be read from the plot by analysing the large error bars. Intermediate $n_\text{hit}$ clusters have the lowest charge-to-energy ratio, as they are more likely to be fully contained. With increasing $n_\text{hit}$, the comp`leteness decreases, leading to a higher charge-to-energy ratio. To assess the impact of this calibration, with the goal of characterizing signal clusters in the context of background with limited the clustering efficiency, one can also investigate a ``cheated charge''. This is computed for all clusters, adding the surrounding charge that has been backtracked to the main electron of the interaction. A calibration based on this cheated clustering gives us an understanding of the losses that arise from a limited clustering ability.

\begin{figure}[ht!]
    \centering
    \includegraphics[width=.49\linewidth]{Images/Reco/charge_correction_factor_hd_1x2x6_marley_official_values_factor_error_fit.png}
    \includegraphics[width=.49\linewidth]{Images/Reco/electroncharge_correction_factor_hd_1x2x6_marley_official_values_factor_error_fit.png}
    \caption{Per hit correction factor distributions. (left) Primary cluster charge. (right) Cheated electron cluster charge.\label{fig:CHARGE_CORRECTION_FACTOR}}
\end{figure}

This trend is well-described by an exponential plus logistic function. The functional form used to fit the charge correction factor as a function of the number of hits is given by:

\begin{align}
   F_0(n_{\text{hit}}) = a \cdot e^{-b \cdot n_{\text{hit}}} + \dfrac{c}{1 + e^{-d \cdot n_{\text{hit}}}}\,,
\end{align}

where \( a, b, c, d \) are fit's exponential amplitude, decay, limiting constant, and sigmoid growth parameters, respectively. This function captures the initial rapid decrease in the correction factor for low hit counts, followed by a more gradual increase at higher hit counts, reflecting the underlying physics of charge collection and clustering in the detector.
\subsection{Electron Energy Calibration}\label{sec:ELECTRON_CALIBRATION}
Even after lifetime and charge calibration, small systematic biases may remain in the reconstructed electron energy due to recombination modelling, imperfect response, or clustering inefficiencies. We apply a linear correction to account for this residual offset:

\begin{equation}
    E_{e}^{\text{final}} = \left(E_{\text{cluster}} - n\right) / m,
\end{equation}

where \( m \) and \( n \) are derived from the linear fits to the simulated \( E_{\text{true}} \) versus \( E_{\text{reco}} \) distributions as shown in \cref{fig:CHARGE_CALIBRATION}. The calibration is performed separately for the main electron cluster and the cheated electron cluster. The calibration factor is computed as a function of the number of hits in the cluster (see~\cref{fig:ENERGY_CALIBRATION}), which serves as a proxy for the electron energy.

\begin{figure}[ht!]
    \centering
    \includegraphics[width=\linewidth]{Images/hd_1x2x6_centralapa_marley_official_primaryenergy_electron_calibration_calibrated_nhits2_hist2d.png}
    \caption[width=1\linewidth]{Charge calibration for the main electron clusters. (left) Comparison between true and reconstructed energy before linear fit calibration. (right) Comparison between true and reconstructed energy after linear fit calibration. Dashed line shows the values for which $E_\text{reco} = E_\text{true}$ as reference.\label{fig:CHARGE_CALIBRATION}}
\end{figure}

As shown in \cref{fig:CHARGE_CALIBRATION_PARAMETERS}, the calibrated parameters approach unity for \( m \) and continuously rise for \( n \) as the number of hits increases, indicating worse electron clustering for higher--energy particles. There is no significant difference between the \gls{hd} and \gls{vd} geometries in terms of these calibration parameters.

\begin{figure}[ht!]
    \centering
    \includegraphics[width=0.49\linewidth]{Images/hd_1x2x6_marley_official_variable_scan.png}
    \includegraphics[width=0.49\linewidth]{Images/vd_1x8x14_3view_30deg_marley_official_variable_scan.png}
    \caption[width=1\linewidth]{Calibration parameters as a function of number of hits in the cluster for (left) \gls{hd} and (right) \gls{vd}.\label{fig:CHARGE_CALIBRATION_PARAMETERS}}
\end{figure}
\subsection{Adjacent Cluster Calibration}\label{sec:GAMMA_CALIBRATION}
Applying this calibration to the main electron cluster, we obtain the final calibrated energy of the main electron. For all the adjacent clusters, we simply apply the average charge correction factor obtained from the main electron calibration.

\begin{figure}[ht!]
    \centering
    \includegraphics[width=.49\linewidth]{Images/hd_1x2x6_marley_official_gamma_energy_iterableNone_hist1d.png}
    \includegraphics[width=.49\linewidth]{Images/hd_1x2x6_centralapa_marley_official_gamma_energy_iterableNone_hist1d.png}
    \caption[width=1\linewidth]{Adjacent cluster signal (left) compared to the selected adjacent clusters (right) for the \gls{hd} productions with and without background activity.\label{fig:RECO_GAMMA_ENERGY_BACKGROUND}}
\end{figure}

We observe a bias in the reconstructed gamma energy (\SI{>2}{\mega\electronvolt}), which is due to the clustering inefficiencies and the low energy of the gammas. This bias leads to a poor identification of the nuclear transitions in the neutrino energy spectrum. In \cref{fig:RECO_GAMMA_ENERGY_BACKGROUND}, we show the reconstructed gamma energy distributions for the different nuclear transition types: Gamow-Teller and Fermi. The Fermi transition has a single emission line at \SI{4.38}{\mega\electronvolt}, while the Gamow-Teller has several and its distribution is therefore divided into two regions: the low-energy region (\SI{<4.38}{\mega\electronvolt}) and the high-energy region (\SI{>4.38}{\mega\electronvolt}).

The bias in the reconstructed gamma energy could be corrected by applying a linear calibration, similar to the one applied to the main electron cluster. However, due to the additional presence of background, it has been considered to focus on improving the channel selection and clustering efficiency without further bias. Additional work needs to be performed in order to estimate the reconstruction potential that can be obtained from correctly tagging these nuclear transitions.

When considering the presence of background activity, as shown in \cref{fig:RECO_GAMMA_ENERGY_BACKGROUND}, the identification of the nuclear transitions becomes even more challenging. The background activity introduces additional low-energy deposits that can be misidentified as gammas from the neutrino interaction, further smearing the reconstructed gamma energy spectrum. This highlights the importance of effective background mitigation strategies and advanced reconstruction techniques to accurately identify and reconstruct the gamma energies associated with neutrino interactions in a realistic detector environment.

\begin{figure}[ht!]
    \centering
    \includegraphics[width=.49\linewidth]{Images/Reco/high_adjacent_cluster_distributions_marley_official_values_counts_adjclenergy_signal_logx_logy_comparison.png}
    \includegraphics[width=.49\linewidth]{Images/Reco/low_adjacent_cluster_distributions_marley_official_values_counts_adjclenergy_signal_logx_logy_comparison.png}
    \caption[width=1\linewidth]{Adjacent cluster distributions for signal and background samples. (left) High-background geometries (\seqsplit{hd\_1x2x6\_lateralAPA} \& \seqsplit{vd\_1x8x14\_3view\_30deg\_nominal}). (right) Low-background geometries (\seqsplit{hd\_1x2x6\_centralAPA} \& \seqsplit{vd\_1x8x14\_3view\_30deg\_shielded}).\label{fig:ADJACENT_CLUSTER_ENERGY_DISTRIBUTION}}
\end{figure}

Another way to illustrate the challenge of reconstructing the gamma energies in the presence of background is to look at the true energy distributions of the adjacent clusters for different background scenarios. \Cref{fig:ADJACENT_CLUSTER_ENERGY_DISTRIBUTION} shows how the presence of background activity significantly increases the number of low-energy adjacent clusters (\SI{<3}{\mega\electronvolt}), which can be misidentified as gammas from the neutrino interaction. From this comparison it is clear that the shielded background scenario reduces considerably the amount of low--energy deposits, but this effect becomes weaker at higher energies.

\subsection{Neutrino Energy Calibration}\label{sec:NEUTRINO_CALIBRATION}
The main electron energy is corrected for the drift time attenuation and the charge calibration factor, as described in \cref{sec:ELECTRON_CALIBRATION}. The main electron energy is then used to compute the total neutrino energy, which is computed from the main electron energy and the energies of all adjacent clusters (up to \SI{1}{\meter}). The adjacent clusters are assumed to be gamma-induced ``blips'' and are calibrated using the same charge calibration factor as the main electron. The total neutrino energy is computed as:

\begin{equation}
    E_{\nu}^{\text{raw}} = E_{e}^{\text{final}} + \sum_i E_{\gamma,i}^{\text{final}}.
\end{equation}

This strategy is the first and most obvious approach for neutrino energy reconstruction (\textbf{TotalEnergy}). However, it does not account for the energy smearing of the signal, and the presence of background can significantly affect the reconstructed neutrino energy. Therefore, we study alternative methods to improve the energy resolution and mitigate the effects of background contamination.

\begin{figure}[ht!]
    \centering
    \includegraphics[width=\linewidth]{Images/hd_1x2x6_centralapa_marley_official_adjcl_selection_signal_hist2d.png}
    \caption[width=1\linewidth]{Distribution of the adjacent clusters charge and distance to the main electron cluster.\label{fig:ADJACENT_CLUSTER_RADIUS_CHARGE}}
\end{figure}

The next simple solution is to apply a selection on the adjacent clusters (\textbf{SelectedEnergy}), which are assumed to be gamma-induced ``blips''. The selection is based on the adjacent cluster charge and distance to the main electron cluster. The adjacent clusters are selected if they have a charge above a certain threshold and are within a certain distance to the main electron cluster. The thresholds are tuned to maximize the signal over background ratio. \cref{fig:ADJACENT_CLUSTER_RADIUS_CHARGE} shows the distribution of the adjacent clusters charge and distance to the main electron cluster and selection thresholds.

\iffalse
\begin{figure}[ht!]
    \centering
    \includegraphics[trim={6cm 0cm 6cm 0cm},clip,width=\linewidth]{Images/hd_1x2x6_centralAPA_marley_official_Random_Forest_Score.png}
    \caption[width=1\linewidth]{(left) Purity vs Efficiency curve of the random forest regressor for both the train and the test samples. (center) Score of the random forest regressor for both \gls{cc} samples. (right) ROC curve of the random forest regressor.\label{fig:RF_SCORE}}
\end{figure}
\fi

\begin{table}[ht!]
    \centering
    \begin{tabular}{lcc}
        \hline
        \textbf{Configuration} & \textbf{Charge (ADC x tick)} & \textbf{Distance (cm)} \\ \hline
        \hline
        hd\_1x2x6 & 5 & 50 \\
        hd\_1x2x6\_centralAPA & 10 & 30 \\
        hd\_1x2x6\_lateralAPA & 8 & 40 \\
        vd\_1x8x14\_3view\_30deg & 12 & 25 \\
        vd\_1x8x14\_3view\_30deg\_nominal & 15 & 20 \\
        vd\_1x8x14\_3view\_30deg\_shielded & 10 & 30 \\
        \hline
    \end{tabular}
    \caption{Adjacent cluster selection thresholds for different background scenarios.\label{tab:ADJACENT_CLUSTER_SELECTION_THRESHOLDS}}
\end{table}

\Cref{tab:ADJACENT_CLUSTER_SELECTION_THRESHOLDS} summarizes the different thresholds used for the adjacent cluster selection, based on the background scenario. The selected adjacent clusters are then summed to the primary cluster energy to obtain the reconstructed neutrino energy (\textbf{SelectedEnergy}).

The next step is to apply a machine learning algorithm to the adjacent clusters to improve the energy resolution. The algorithm is trained on the main cluster properties and the average adjacent cluster features. The algorithm is a random forest regressor, which is trained on the simulated data. The random forest regressor is used to predict the energy of the adjacent clusters, which are then added to the main electron energy to obtain the total neutrino energy (\textbf{SolarEnergy}). To simply the problem and motivated by the two main calibration lines from the \gls{cc} interaction displayed in~\cref{fig:MARLEY_CC_INTERACTION}, the \gls{ml} algorithm is trained two add one of two total adjacent cluster energy values: \qtylist{4.47;6.41}{\mega\electronvolt}.\iffalse \cref{fig:RF_SCORE} shows \fi The score of the random forest regressor which provides an AUC of \num{0.77} for the test sample, which means that the algorithm is able to distinguish between signal and background with a good efficiency. The random forest regressor is trained on the following features:

\begin{itemize}
    \item NHits: number of hits in the main electron cluster.
    \item AdjClNum: number of total adjacent clusters.
    \item SelectedAdjClNum: number of selected adjacent clusters.
    \item Energy: energy of the main clusters.
    \item \textbf{TotalAdjClEnergy}: energy of the total adjacent clusters.
    \item SelectedAdjClEnergy: energy of the selected adjacent clusters.
    \item SelectedAdjClEnergyRatio: energy ratio of the selected to the total adjacent clusters.
    \item \textbf{MaxAdjClEnergy}: maximum energy of the total adjacent clusters.
    \item SelectedMaxAdjClEnergy: maximum energy of the selected adjacent clusters.
    \item TotalEnergy: reconstructed \textbf{TotalEnergy}.
    \item SelectedEnergy: reconstructed \textbf{SelectedEnergy}.
\end{itemize}

From which the highlighted features are the most important for the random forest regressor, as shown in \cref{fig:RF_FEATURE_IMPORTANCE}. Finally, as a benchmark for comparison between the here presented reconstruction algorithms and with the aim of quantizing the impact of adjacent gammas on the total neutrino energy reconstruction, we introduce the \textbf{ClusterEnergy} algorithm as the reconstruction method based purely on the main electron cluster. The only correction is a fixed calibration offset applied to the full sample to align the center of the distribution with the true energy.

\begin{figure}[ht!]
    \centering
    \includegraphics[width=.49\linewidth]{Images/hd_1x2x6_centralAPA_marley_official_Random_Forest_Feature_Importance.png}
    \includegraphics[width=.49\linewidth]{Images/hd_1x2x6_centralAPA_marley_official_Discriminant_Calibration_1D.png}
    \caption[width=1\linewidth]{(left) Feature importance of the random forest regressor. (right) Discriminant calibration for the random forest regressor.\label{fig:RF_FEATURE_IMPORTANCE}}
\end{figure}


\subsection{Neutrino Energy Smearing\label{sec:SIGNAL_SMEARING}}
Due to the missing nuclear binding energy and undetected soft particles, the reconstructed neutrino energy is expected to systematically underestimate the true value. As a final step to the reconstruction process, all energies (\textbf{TotalEnergy}, \textbf{SelectedEnergy}, \textbf{SolarEnergy} and \textbf{ClusterEnergy}) are calibrated with a linear fit to compensate for systematic losses in the previous clustering or calibration processes. This final global correction is applied to bring the distribution into agreement with the simulation truth:

\begin{equation}
    E_{\nu}^{\text{cal}} = \left(E_{\nu}^{\text{raw}} - n\right) / m,
\end{equation}

with \( m \) and \( n \) obtained from a linear fit to the true versus reconstructed neutrino energy spectrum of the \gls{mc} sample as shown in \cref{fig:NEUTRINO_ENERGY_SMEARING}.

\begin{figure}[ht]
    \centering
    \includegraphics[width=\linewidth]{Images/hd_1x2x6_centralapa_marley_official_neutrino_energy_variable_hist2d.png}
    \caption[width=1\linewidth]{Neutrino energy smearing from the calibrated signal.\label{fig:NEUTRINO_ENERGY_SMEARING}}
\end{figure}